% Part: sets-functions-relations
% Chapter: functions
% Section: basics

\documentclass[../../../include/open-logic-section]{subfiles}

\begin{document}

\olfileid[fa-IR]{sfr}{fun}{bas}
\olsection{مبانی}

\begin{explain}
یک \emph{تابع} نگاشتی است که هر !!{element} از یک مجموعهٔ معین را
به یک !!{element} معین در یک مجموعهٔ معینِ (دیگر) می‌فرستد. برای نمونه،
عمل افزودن~$1$ تابعی را تعریف می‌کند: هر عدد~$n$ به
عدد یکتای~$n+1$ نگاشته می‌شود.

به‌طور کلی‌تر، توابع ممکن است زوج‌ها، سه‌تایی‌ها و مانند آن‌ها را ورودی بگیرند و
نوعی خروجی بازگردانند. بسیاری از توابع را از
حساب مقدماتی می‌شناسیم. برای نمونه، جمع و ضرب
تابع‌اند. آن‌ها دو عدد را ورودی می‌گیرند و عدد سومی بازمی‌گردانند.

در این معنای ریاضی و انتزاعی، تابع یک \emph{جعبهٔ
سیاه} است: تنها این مهم است که کدام خروجی با کدام ورودی جفت می‌شود، نه
روش محاسبهٔ خروجی.
\end{explain}

\begin{defn}[تابع]
یک \emph{تابع} $f \colon A \to B$ نگاشتی از هر !!{element}
از~$A$ به یک !!{element} از~$B$ است.

$A$ را \emph{دامنهٔ}~$f$ و $B$ را \emph{هم‌دامنهٔ}
~$f$ می‌نامیم. !!{element}s از~$A$، ورودی‌ها یا \emph{آرگومان‌های}
~$f$ نامیده می‌شوند، و !!{element} از~$B$ که با آرگومان~$x$
به‌وسیلهٔ~$f$ جفت می‌شود، \emph{مقدار~$f$} برای آرگومان~$x$ نام دارد و
به‌صورت~$f(x)$ نوشته می‌شود.

\emph{برد} $\ran{f}$ تابع~$f$ زیرمجموعه‌ای از هم‌دامنه است
که از مقادیر~$f$ برای برخی آرگومان‌ها تشکیل می‌شود؛ $\ran{f} =
\Setabs{f(x)}{x \in A}$.
\end{defn}

نمودارِ \olref{fig:function} می‌تواند به فهم توابع کمک کند. بیضی
سمت چپ \emph{دامنهٔ} تابع را بازنمایی می‌کند؛ بیضی سمت
راست \emph{هم‌دامنهٔ} تابع را بازنمایی می‌کند؛ و پیکانی
از یک \emph{آرگومان} در دامنه به
\emph{مقدار} متناظر در هم‌دامنه اشاره می‌کند.

\begin{figure}
  \olasset{assets/diagrams/function.tikz}
  \caption{یک تابع نگاشتی از هر !!{element} از یک مجموعه به
    !!a{element} از مجموعه‌ای دیگر است. پیکانی از یک آرگومان در
    دامنه به مقدار متناظر در هم‌دامنه اشاره می‌کند.}
  \ollabel{fig:function}
\end{figure}

\begin{ex}
ضرب، زوج‌های اعداد طبیعی را ورودی می‌گیرد و آن‌ها را
به اعداد طبیعی به‌عنوان خروجی نگاشت می‌کند؛ پس از $\Nat \times \Nat$ (
دامنه) به $\Nat$ (هم‌دامنه) می‌رود. چنان‌که معلوم می‌شود، برد نیز
$\Nat$ است، زیرا هر $n \in \Nat$ برابر $n \times 1$ است.
\end{ex}

\begin{ex}
ضرب یک تابع است، زیرا هر ورودی---هر زوج
از اعداد طبیعی---را با یک خروجی یگانه جفت می‌کند: $\times \colon \Nat^2 \to
\Nat$. در مقابل، نگاشت‌کردن عدد طبیعی~$n$ به عدد حقیقی~$x$ چنان‌که
$x^2 = n$، تابع‌وار نیست، زیرا هر عدد صحیح مثبت~$n$ دو
ریشهٔ دوم دارد: $\sqrt{n}$ و~$-\sqrt{n}$. می‌توانیم با
بازگرداندن تنها ریشهٔ دوم مثبت آن را تابع‌وار کنیم: $\sqrt{\phantom{X}} \colon
\Nat \to \Real$.
\end{ex}

\begin{ex}
رابطه‌ای که هر دانش‌آموز یک کلاس را با نمرهٔ نهایی او جفت می‌کند،
تابع است---هیچ دانش‌آموزی نمی‌تواند در همان
کلاس دو نمرهٔ نهایی متفاوت بگیرد. رابطه‌ای که هر دانش‌آموز یک کلاس را با
والدین او جفت می‌کند تابع نیست: دانش‌آموزان می‌توانند صفر، دو، یا بیش از
دو والد داشته باشند.
\end{ex}

\begin{explain}
می‌توانیم توابع را با تعیین دقیق مقدار تابع برای
هر آرگومان ممکن تعریف کنیم. روش‌های گوناگون این کار عبارت‌اند از
ارائهٔ یک فرمول، توصیف روشی برای محاسبهٔ
مقدار، یا فهرست‌کردن مقادیر برای هر آرگومان. توابع را به هر شیوه‌ای که
تعریف کنیم، باید اطمینان یابیم که برای هر آرگومان یک،
و تنها یک، مقدار تعیین می‌کنیم.
\end{explain}


\begin{ex}
فرض کنید $f \colon \Nat \to \Nat$ چنان تعریف شود که $f(x) = x+1$. این
تعریف، $f$ را به‌صورت تابعی مشخص می‌کند که
اعداد طبیعی را ورودی می‌گیرد و اعداد طبیعی خروجی می‌دهد. این تعریف می‌گوید که با داشتن یک
عدد طبیعی~$x$، تابع $f$ جانشین آن، یعنی~$x+1$، را خروجی خواهد داد.
در این مورد، هم‌دامنهٔ $\Nat$ بردِ~$f$ نیست، زیرا
عدد طبیعی~$0$ جانشین هیچ عدد طبیعی‌ای نیست.
بردِ~$f$ مجموعهٔ همهٔ اعداد صحیح مثبت، یعنی $\Int^{+}$، است.
\end{ex}

\begin{ex}\ollabel{examplefunext}
فرض کنید $g \colon \Nat \to \Nat$ چنان تعریف شود که $g(x) = x+2-1$. این
تعریف به ما می‌گوید که $g$ تابعی است که اعداد طبیعی را ورودی می‌گیرد و
اعداد طبیعی خروجی می‌دهد. با داشتن یک عدد طبیعی~$n$، تابع $g$
پیشینِ جانشینِ جانشینِ~$x$، یعنی
$x+1$، را خروجی خواهد داد.
\end{ex}

\begin{explain}
اکنون دو تابع $f$ و $g$ را با \emph{تعریف‌های} متفاوت بررسی کردیم.
بااین‌حال، این‌ها \emph{همان تابع} هستند. چراکه
برای هر عدد طبیعی~$n$ داریم $f(n) = n+1 = n+2-1 =
g(n)$. به بیان دیگر، تعریف‌های ما برای $f$ و~$g$ یک
نگاشت واحد را به‌وسیلهٔ معادلات متفاوت مشخص می‌کنند. پس به‌طور ضمنی
بر اصل گسترش برای توابع تکیه می‌کنیم،
\[
  \text{اگر }\forall x\, f(x) = g(x)\text{، آنگاه }f = g
\]
مشروط بر اینکه $f$ و~$g$ دامنه و هم‌دامنهٔ یکسانی داشته باشند.
\end{explain}

\begin{ex}
توابع را می‌توانیم به‌صورت موردی نیز تعریف کنیم. برای نمونه، می‌توانیم
$h \colon \Nat \to \Nat$ را چنین تعریف کنیم:
\[
h(x) =
\begin{cases}
  \frac{x}{2} & \text{اگر $x$ زوج باشد} \\
  \frac{x+1}{2} & \text{اگر $x$ فرد باشد.}
\end{cases}
\]
چون هر عدد طبیعی یا زوج است یا فرد، خروجی این
تابع همواره عددی طبیعی خواهد بود. فقط به یاد داشته باشید که اگر
تابعی را به‌صورت موردی تعریف می‌کنید، هر ورودی ممکن باید در
دقیقاً یک حالت قرار گیرد. در برخی موارد، این امر مستلزم اثبات آن است که
حالت‌ها جامع و دوبه‌دو ناسازگارند.
\end{ex}

\end{document}
